\section{Accuracy versus quantization versus context depth}\label{sec:depth}

Does a 4-bit build lose more than a 6-bit build as the context fills? Divergence cannot be measured at
depth on this host (\cref{sec:hardware}), and comparing free-running generations fails for a structural
reason: I tried it, and greedy trajectories from different builds fork within the first 2 to 126
characters, after which every pair of builds is ``completely different'' (normalised edit distance
0.61--0.79) regardless of how close they are on the ladder. Adjacent builds measured further apart than
the extremes. Distance between quantizations has to be measured on distributions at a fixed context,
not on emitted text. That leaves task instruments: RULER in the first battery, AA-LCR and GPQA in the
second.

\subsection{RULER retrieval to 131K, and a headline withdrawn}\label{sec:ruler}

\begin{table}[t]
\centering\small
\caption{RULER, first battery: RULER's own generators and \texttt{string\_match\_all} metric, greedy,
no speculation, one split for both builds, \flag{q4\_0} KV. The harness used the raw completion endpoint
with a \textbf{128-token output budget and reasoning left on}. ``Closed'' means the output contains
the end-of-reasoning tag. Prompts reached 98--100\,\% of the target length in every cell.}\label{tab:ruler}
\shrinktofit{\begin{tabular}{@{}lrlrrrr@{}}
\toprule
task & context & build & $n$ & score\,\% & closed & closed \emph{and} wrong \\
\midrule
single needle & 8K / 32K / 131K & \QsixXL & 25 / 25 / 12 & 100 / 100 / 100 & all & 0 \\
single needle & 8K / 32K / 131K & \Qfour & 25 / 25 / 12 & 100 / 100 / 100 & all & 0 \\
multi-key (4 keys) & 131{,}072 & \QsixXL & 100 & \textbf{89.0} & 77 & \textbf{0} \\
multi-key (4 keys) & 131{,}072 & \Qfour & 100 & \textbf{79.0} & 60 & \textbf{0} \\
\midrule
\multicolumn{7}{@{}l}{multi-key, paired: score discordance 10 vs 0, exact McNemar $p=0.002$, $-10.0$ pts $[-15.9,-4.1]$} \\
\multicolumn{7}{@{}l}{\phantom{multi-key, paired: }closure discordance 22 vs 5, $p=0.0015$, $-17.0$ pts $[-26.6,-7.4]$} \\
\multicolumn{7}{@{}l}{\phantom{multi-key, paired: }on the 55 items where \emph{neither} build hit the budget: 55/55 vs 55/55, zero discordance} \\
\bottomrule
\end{tabular}}
\end{table}

\textbf{Single-needle retrieval is saturated and unaffected.} Both ends of the ladder score 100\,\% at
8K, 32K and 131K tokens, and every response finished its reasoning inside the budget, so these are
real ceilings, not artifacts of lenient string matching (\cref{tab:ruler}).

\textbf{The multi-key result is a budget effect.} With four keys planted and one queried, \QsixXL{}
scored 89/100 and \Qfour{} 79/100 at 131K tokens. All ten discordant items favoured the larger build;
$p=0.002$; recovery 88.8\,\%, right inside the band Red Hat publishes for 4-bit weights at 128K
\citep{redhat_longcontext_quant}. I wrote it up as the study's long-context headline. It took 9.4 GPU-hours and
the analysis script was committed before the second arm finished.

It does not measure retrieval. Re-reading the saved predictions: \emph{not one response in either arm
finished reasoning and then gave a wrong answer.} Every failure is a truncation at 128 tokens. On the
55 items where neither build hit the budget, both score 55 of 55. What differs is how often the
reasoning block closed in time: 77 against 60, itself significant ($p=0.0015$). The 4-bit build
``thought'' longer on the harder task. The closure audit also shows why depth is not the cause: at the
same 131K tokens, same builds, same budget, single-needle closes 12 of 12 and multi-key closes 77 and
60 of 100. Task difficulty drives the reasoning length.

So the retrieval question at 131K on the multi-key task is \emph{unanswered}, not answered in the
negative, and the comparison with the Red Hat band is void. What remains is a smaller finding of its own:
a more heavily quantized build spent measurably more tokens reasoning about the same questions. The
same budget mechanism explains the thinking-mode HumanEval+ rows in \cref{tab:humaneval}, and a harder
RULER task (variable tracking) was excluded from the battery for the same defect in its extreme form
(0 of 25 closed). A score alone cannot tell a real ceiling from a truncation; the closure column can.

\subsection{AA-LCR and GPQA Diamond: long-document reasoning with a short-context control}\label{sec:aalcr}

The second battery avoided those traps. It used the chat endpoint with the model's documented
reasoning control at its highest effort (\texttt{xhigh}), Qwen's official thinking sampling, and an
output budget large enough that truncation was rare (1 of 756 answers). Two instruments ran under
identical settings:

\textbf{AA-LCR} \citep{artificialanalysis2026aalcr}: 100 questions, each over a set of real documents
totalling about 100K tokens, requiring reasoning across them. Served at the full 262K window with each
build's gentlest workable cache: \flag{q4\_0} for \Qsix, \flag{q8\_0} for \Qfour. Answers took 2 to 60
minutes each. Correctness is judged by an LLM equality checker using the benchmark's published prompt.

\textbf{GPQA Diamond} \citep{rein2023gpqa}: 198 graduate-level science questions of about 200 tokens,
scored by answer extraction. This is the short-context control: if a drafter or build broke reasoning
in general, it would show here. Single answers reached 94{,}000 generated tokens. With short prompts
the cache could be gentler still: \flag{q8\_0} for \Qsix{} and \flag{f16} for \Qfour.

Design: all four build $\times$ drafter configurations ran a paired subset (30 AA-LCR questions, 50
GPQA items); then the faster drafter (DFlash2, \cref{sec:s15speed}) ran the full sets on both builds.
A full pass with MTP was not run, by design, to stay inside a 117-hour GPU budget. Build and cache
precision change together between arms, so the comparison is between \emph{configurations}, never
between quantizations alone.

\begin{figure}[t]
\centering
\begin{tikzpicture}
\begin{axis}[paper,height=7.6cm,width=0.8\linewidth,
  xmin=50,xmax=100,xlabel={accuracy (\%), Wilson 95\,\% interval},
  ytick={1,2,3,4,5.5,6.5,8,9,10,11,12.5,13.5},
  yticklabels={{\Qsix{} + MTP, $n{=}30$},{\Qsix{} + DFlash2, $n{=}30$},{\Qfour{} + MTP, $n{=}30$},{\Qfour{} + DFlash2, $n{=}30$},{\Qsix{} + DFlash2, $n{=}100$},{\Qfour{} + DFlash2, $n{=}100$},{\Qsix{} + MTP, $n{=}50$},{\Qsix{} + DFlash2, $n{=}50$},{\Qfour{} + MTP, $n{=}50$},{\Qfour{} + DFlash2, $n{=}50$},{\Qsix{} + DFlash2, $n{=}198$},{\Qfour{} + DFlash2, $n{=}198$}},
  yticklabel style={font=\scriptsize},ymin=0.2,ymax=14.9,
  error bars/x dir=both,error bars/x explicit,error bars/error bar style={line width=0.8pt}]
\addplot[only marks,mark=*,cbblue] coordinates {
 (73.3,1) += (12.5,0) -= (17.8,0)
 (80.0,2) += (10.5,0) -= (17.3,0)
 (83.3,3) += (9.3,0) -= (16.9,0)
 (86.7,4) += (8.0,0) -= (16.4,0)
 (83.0,5.5) += (6.1,0) -= (8.5,0)
 (87.0,6.5) += (5.2,0) -= (8.0,0)};
\addplot[only marks,mark=square*,cbred] coordinates {
 (96.0,8) += (2.9,0) -= (9.5,0)
 (94.0,9) += (3.9,0) -= (10.2,0)
 (94.0,10) += (3.9,0) -= (10.2,0)
 (90.0,11) += (5.7,0) -= (11.4,0)
 (89.9,12.5) += (3.5,0) -= (5.0,0)
 (88.4,13.5) += (3.7,0) -= (5.2,0)};
\draw[black!30] (axis cs:50,7.25) -- (axis cs:100,7.25);
\node[font=\footnotesize\bfseries,text=cbblue,anchor=west] at (axis cs:51,6.7) {AA-LCR ($\sim$100K-token documents)};
\node[font=\footnotesize\bfseries,text=cbred,anchor=west] at (axis cs:51,14.35) {GPQA Diamond (short prompts)};
\draw[dashed,cbblue] (axis cs:82.0,0.2) -- (axis cs:82.0,7.25);
\draw[dashed,cbred] (axis cs:90.5,7.25) -- (axis cs:90.5,14.9);
\end{axis}
\end{tikzpicture}
\caption{No build and no drafter separates from any other on either instrument. Second battery, build
E-C, reasoning effort \texttt{xhigh}, official thinking sampling. AA-LCR rows are the three-judge
resolved scores; the $n{=}30$ DFlash2 rows are scored on the paired-subset items, where the resolved
and first-judge scores coincide. Dashed lines are Artificial Analysis' figures for the same model served through the vendor's API
(0.820, 0.905), shown for orientation: different serving stack, precision and, for AA-LCR, judge. Build and cache precision change together between rows: the comparison is between configurations, never between quantizations alone.}\label{fig:accuracy}
\end{figure}

\begin{table}[t]
\centering\small
\caption{AA-LCR judging. The same 260 answers were scored by two independent LLM judges using the
benchmark's equality-checker prompt; a third judge ruled on the 11 disagreements. Identifiers are as
recorded in the artifacts.}\label{tab:judges}
\shrinktofit{\begin{tabular}{@{}lrrrrr@{}}
\toprule
judge & \Qfour+DFlash2 & \Qfour+MTP & \Qsix+DFlash2 & \Qsix+MTP & total \\
 & (/100) & (/30) & (/100) & (/30) & (/260) \\
\midrule
judge 1 (\texttt{muse-spark}, effort xhigh) & 88 & 25 & 84 & 22 & 219 \\
judge 2 (\texttt{gpt-5.6-luna}, effort medium) & 87 & 25 & 87 & 23 & 222 \\
\textbf{resolved} (third judge on disagreements) & \textbf{87} & \textbf{25} & \textbf{83} & \textbf{22} & \textbf{217} \\
\midrule
\multicolumn{6}{@{}l}{Cohen's $\kappa$ over all 260: judge 1 $\times$ 2: 0.836 (agreement 249/260); 1 $\times$ 3: 0.972; 2 $\times$ 3: 0.869} \\
\bottomrule
\end{tabular}}
\end{table}

\paragraph{Result.} \Cref{fig:accuracy}: every interval overlaps every other within each instrument.
On the full AA-LCR pass \Qfour{} scores 87/100 and \Qsix{} 83/100 (resolved); on full GPQA, 175/198 and
178/198. The nominal order flips between the two instruments, as it did between HumanEval+ and
SWE-bench in the first period. On the paired subsets the drafters do not separate either (AA-LCR
22--26 of 30; GPQA 45--48 of 50). Speculation at depth 7 on long reasoning chains did not cost
accuracy that these instruments can see, at $n=30$ to 198.

This is the second battery's version of the report's thesis: 756 long answers, about 69 GPU-hours,
and a 4-bit build with an 8-bit cache is indistinguishable from a 6-bit build with a 4-bit cache at
the full window. Divergence ordered the same two weight builds in every one of 32 windows, in forty minutes.
Both statements are true. They answer different questions, and only the first is about outcomes a
user would notice.

\paragraph{Judging.} An LLM-judged benchmark carries judge variance, and a two-of-three majority is a
consensus label, not ground truth. The Wilson intervals in \cref{fig:accuracy} cover neither judge error
nor the clustering of questions within shared document sets. \Cref{tab:judges}: the two full
judges differ by 3 items overall and by at most 3 on any arm; the resolved score sits in a 5-item band
(217 to 222). A hand audit of the disagreements found no pipeline defect. It did find that the misses
are not scattered: a small set of questions recurs across arms where the official key is arguable
(a rounding convention, a conditional answer scored against an unconditional key, an entity list with
one debatable member). Four to six items decide every subset contrast, and several of those are
key-ambiguous. AA-LCR scores should always name the judge and its settings.

\paragraph{Against the public baseline.} Artificial Analysis reports 0.820 on AA-LCR and 0.905 on GPQA
Diamond for this model at \texttt{xhigh} effort through Alibaba's API
\citep{artificialanalysis2026qwen38}. The local figures (83--87\,\% and 88--90\,\%) are in the same
range. The comparison is loose in several ways: a hosted model presumably at full precision against
local 4- to 6-bit builds with quantized caches; an undisclosed sampling temperature on their side;
for AA-LCR a different judge. It supports ``local quantized serving is in the neighbourhood of the
hosted model on these two instruments'', and no more.

\subsection{What depth does to speed}
Depth costs speed far more visibly than it costs accuracy. Without speculation, decode falls from
11.7--13.1\tps{} at 32K filled tokens to 3.5--3.6 at 246K (\cref{fig:specdepth}): attention reads the
whole cache for every token, and under layer split the card that owns a layer does all of that work.
Prefill falls from about 490--530\tps{} to 215--230. With DFlash2 the same 246K context still decodes
at 26--28\tps. A speed table measured into an empty cache (which is what most published tables are,
including my own early ones at 57.5\tps{}) overstates full-window speed by $3$--$5\times$. Any speed
figure should state how full the cache was.

\begin{takeaway}
Up to 131K tokens, single-needle retrieval is perfect at 4 and 6 bits. At about 100K tokens of real
documents, reasoning accuracy is 83--87\,\% on the two full passes and 73--87\,\% on the 30-question
subsets, for both builds and both drafters, with overlapping intervals and a nominal order that flips between instruments. No accuracy instrument in this study
shows a quantization $\times$ depth interaction; none had the power to exclude a small one. Output
budget, not retrieval, produced the one ``significant'' long-context difference.
\end{takeaway}
